\documentclass{article}
\usepackage{amsmath,amssymb,amsthm,hyperref}
\newtheorem{theorem}{Theorem}\newtheorem{corollary}[theorem]{Corollary}
\title{Native-natural coupling selection forces the independent law}
\date{16 November 2026}
\begin{document}\maketitle
\paragraph{Audit, attribution and declared kernel.}
The scalar uniqueness result below is an application of an established theorem,
not a novelty claim. Bloch and Fritz, \emph{Naturality characterizes product
measures and relative product measures}, arXiv:2609.15909v2, Theorem 3.1,
prove uniqueness for tuples on measurable spaces and explicitly for finite
sets; Theorem 3.2 treats a fixed shared base. The proof supplied here is an
elementary finite-interface derivation adapted to retained raw witnesses.
The native consequence concerns the selected \emph{scalar joint law}; it
neither identifies tensor and product coefficients nor constructs a module,
contentad, substitution compiler or coherent selector of actual coefficient
arrows. No continuity assumption or deletion of zero witnesses is used.

\section{Exact interface}
Let $\mathsf{FinProb}_{\rm raw}$ have objects $(W,p)$ with $W$ a finite set
and $p:W\to[0,1]$, $\sum_wp(w)=1$. Its arrows are all actual functions
$r:W\to W'$ satisfying
\[
 p'(w')=\sum_{r(w)=w'}p(w).
\]
Zero-weight elements remain elements of $W$; arrows act on them too.
No normalized object has empty $W$. Fix a finite arity $n$.
A selection assigns every ordered tuple $((W_i,p_i))_{i=1}^n$ a probability
law $J(p_1,\ldots,p_n)$ on the \emph{full} raw product $\prod_iW_i$.
For $n>0$ require its marginals to be $p_i$. Require, for every independently
chosen tuple of arrows $r_i:(W_i,p_i)\to(W_i',p_i')$,
\[
 (\prod_i r_i)_*J(p_1,\ldots,p_n)=J(p_1',\ldots,p_n'). \tag{N}
\]
This is naturality in the product category, not just simultaneous relabelling
of one common input, not just permutation covariance, and not naturality on
an already selected class of coupling-preserving maps. The selected law
is allowed arbitrary set-theoretic dependence on its arguments before (N)
is imposed. No associativity, symmetry, continuity, measurability in $p_i$,
or monad compatibility is an assumption.

\begin{theorem}[Attributed finite selector uniqueness]
For each finite $n$, the only selection satisfying this interface is
\[
 J(p_1,\ldots,p_n)(w_1,\ldots,w_n)=\prod_{i=1}^np_i(w_i).
\]
For $n=0$ this is the unique mass-one law on the singleton empty product;
for $n=1$ it is the input law.
\end{theorem}
\begin{proof}
The cases $n=0,1$ are immediate; the following proof also works for $n=1$.
Write $B_p$ for the law on the full two-element set $\{0,1\}$ with
mass $p$ at $1$, even at $p=0$ or $1$. Define
\[
 f(p_1,\ldots,p_n)=J(B_{p_1},\ldots,B_{p_n})(1,\ldots,1).
\]
For subsets $S_i\subseteq W_i$, their indicator functions are admitted
probability-pushforward arrows to $B_{p_i(S_i)}$. Naturality therefore gives
\[
 J(p_1,\ldots,p_n)(S_1\times\cdots\times S_n)
 =f(p_1(S_1),\ldots,p_n(S_n)). \tag{R}
\]
All such subsets, including empty and full subsets, are available.

Fix one coordinate and all other Bernoulli probabilities. For $a,b\geq0$
with $a+b\leq1$, put the law $(a,b,1-a-b)$ on the full three-atom raw set
$\{A,B,C\}$ in that coordinate. Use its three indicator maps for
$\{A\}$, $\{B\}$ and $\{A,B\}$, leaving the other coordinates unchanged.
The corresponding product events in the first two cases are disjoint.
Finite additivity of this one selected joint law and (N) give
\[
 f(\ldots,a+b,\ldots)=f(\ldots,a,\ldots)+f(\ldots,b,\ldots). \tag{A}
\]
The argument includes zero masses and requires no positivity assumption.

For the resulting function $g:[0,1]\to[0,1]$, (A) implies $g(0)=0$
and, by nonnegativity, $g$ is monotone: $g(y)-g(x)=g(y-x)\geq0$ when
$x\leq y$. Repeated addition gives $g(k/m)=(k/m)g(1)$.
For arbitrary $x$, rational lower and upper bounds and monotonicity squeeze
$g(x)$ to $xg(1)$. This derives linearity without assuming continuity.
Applying this argument to each coordinate gives
\[
 f(p_1,\ldots,p_n)=\bigl(\prod_i p_i\bigr)f(1,\ldots,1).
\]
The coupling with all marginals $B_1$ is supported on the all-ones tuple,
so $f(1,\ldots,1)=1$. Taking each $S_i$ in (R) to be a singleton proves
the displayed formula on every raw tuple, including zero-weight tuples.
Conversely finite distributivity of sums proves (N) for the product law,
and its marginals are the prescribed laws. Thus existence and uniqueness
both hold.
\end{proof}

\paragraph{The marginal axiom can be omitted.}
The proof remains valid if $J$ is only required to output a normalized
nonnegative law, with no marginal axiom. Indeed the all-singleton input has
only one possible output law. Mapping each singleton to the value $1$ of
$B_1$ forces $f(1,\ldots,1)=1$ by (N); the rest of the proof is unchanged.
Thus the stronger formulation agrees with the form of the literature theorem.
Normalization and nonnegativity of the selected law are retained assumptions.

\section{Countable raw sets and zero witnesses}
The same statement holds if the input category includes countable raw sets
with finitely supported normalized laws and all exact-pushforward functions.
Indicator maps and the finite three-atom tests still belong to that category,
so the proof gives the same weight on every tuple. This retains every raw zero
witness: it specifies zero mass, not removal or identification. Alternatively,
finite supports map by actual inclusions into the original raw sets, and (N)
forces the product law as the pushforward of the finite result. These
inclusions are proof probes, not a change of the output raw interface.

The argument even extends to arbitrary normalized laws on countable raw sets,
provided the selected output is a countably additive discrete probability law.
A finite product of countable sets is countable, and singleton weights determine
its measure. No convergence assumption on the selection is introduced.
This statement is for finite arity; no infinite-product interface is supplied.
For arbitrary finite measures of unequal total mass the coupling interface
itself changes or can be empty, so that setting is not silently included.

\section{Actual augmented coefficients}
Let $\mathcal E$ be a coefficient category with terminal object $\mathbf1$.
A decorated normalized row has $(W,p,(E_w)_{w\in W})$. A native arrow
consists of an exact-pushforward raw map $r$ and actual arrows
$E_w\to E'_{r(w)}$. Suppose a native selection has raw product interface,
selects a scalar coupling, and its comparison over any tuple of native arrows
has raw function $\prod_i r_i$ and preserves the selected scalar law by exact
pushforward. Its coefficient construction can otherwise remain unspecified.

\begin{corollary}[Coefficient-erasure test]
Under these assumptions, the scalar law of any such native-natural selection
is the product law, independently of the decorations.
\end{corollary}
\begin{proof}
There is an actual native arrow from each decorated row to its terminally
decorated version, identity on the entire raw set and with the same weights.
Applying the selection to these arrows gives identity raw map on the entire
product, so exact pushforward forces the two selected scalar laws to agree.
On terminally decorated rows every deterministic probability-pushforward
function lifts to a native arrow via the unique coefficient maps. The theorem
therefore applies on this subcategory, and hence to the original rows.
\end{proof}

For the augmented vector category $\mathrm{FinVect}_k/k$, the terminal object
is $k\xrightarrow{1}k$, and the coefficient-erasure arrow is the actual
augmentation $E_w\to k$. This is not a faithful-carrier argument. Both
tensor-style and product-style selected coefficients can be subject to the
same scalar test when their native comparison data satisfy the stated
hypothesis; the theorem neither equates those coefficients nor asserts that
all required output arrows exist.
If the coefficient category lacks the erasure arrows, the theorem still
applies on any admitted terminal/unit-decorated subcategory carrying all raw
pushforward maps, but it does not show that laws on other decorations agree.
The same caution applies to content labels, fixed outputs and other material
structure: no erasure map for them is assumed. A fixed body may obstruct the
indicator tests. The conclusion applies to the bare row/joint-context
interface above, or to a larger native interface only after the required maps
are explicitly shown to lift.

\section{Sharp boundaries and relevance}
The theorem forbids a control-free, scalar-correlating selection on this
entire interface while preserving all its native maps. It does not forbid
correlated couplings or a family containing every coupling. A specified
coupling is additional data, and its compatible maps form a smaller category.
A controller may also choose joint laws at each retained causal scope; its
comparisons must transport those chosen laws rather than all independent
marginal maps indiscriminately.

For example, for a single common $(W,p)$ the diagonal law on $W\times W$
is natural under synchronous maps $r\times r$. It is not natural under
independent coordinate maps: for a fair bit the first-coordinate flip sends
the diagonal coupling to the antidiagonal coupling, while leaving both input
marginals unchanged. Thus identifying two input occurrences by a common label
or shared latent variable changes the morphism interface materially.
Likewise, working over a retained common base can permit unconditional
correlations: independent arbitrary coordinate maps have been replaced by
maps over that base. The literature's relative theorem is relevant to that
altered interface, not a counterexample to the theorem above.

The relative result fixes both the base and its probability law. Its
conditional distributions are determined only up to base-almost-everywhere
equality, whereas the integrated joint law is uniquely well-defined and
independent of those choices. It supplies neither a complete conditional
control table on zero-probability histories nor arbitrary base-changing
prefix transport. Those distinctions remain material in the causal strategy
interface.

\paragraph{A finite base-coarsening diagnostic.}
Fix a finite probability base $(B,\omega)$ and $g:B\to C$, with
$\nu=g_*\omega$. For a fixed $n\geq2$, let normalized rows $(X_i,p_i)$
carry maps $\rho_i:X_i\to B$ with $(\rho_i)_*p_i=\omega$.
There is an inclusion of full raw pullbacks
\[
 j:\prod_B X_i\longrightarrow\prod_C X_i.
\]
The conditional-product law over $B$ assigns an aligned tuple with base $b$
and $\omega(b)>0$ weight
$\prod_i p_i(x_i)/\omega(b)^{n-1}$, and assigns zero weight otherwise.
Define the law over $C$ by the analogous formula with $\nu$ and $g\rho_i$.
Then $j$ preserves these laws by exact pushforward for \emph{every} such
row family if and only if each fibre of $g$ contains at most one
positive-$\omega$ point.
Indeed, under this condition every positive tuple aligned over $C$ is already
aligned over $B$, its unique positive base weight equals the $C$-base weight,
and the two displayed formulas agree. All other tuples have zero weight.
Conversely choose every $X_i=B$, $p_i=\omega$, $\rho_i=1_B$.
Two distinct positive points $b,b'$ in one $g$-fibre give a positive target
tuple $(b,b',b,\ldots,b)$ outside the source pullback. Exact pushforward
along $j$ is impossible. This is an elementary finite diagnostic, not a claim
of new literature-level uniqueness. Even in the permitted case $j$ need not
be a raw bijection: extra zero-weight tuples in the target remain retained.
Thus scalar compatibility alone does not imply a native isomorphism.
For $n=1$ both laws are the original law and this obstruction disappears.

A simple failure of the three-atom test occurs for a proposed Bernoulli
selector with joint all-ones probability $f(p,q)=\min(p,q)$.
At $q=1/2$, $a=b=3/10$, equation (A) would require
$1/2=3/10+3/10$, which is false. Arbitrary dependence on real-valued weights
does not evade the result: order and finite additivity already force linearity.

The result concerns selecting one joint law, not existence of universal joint
objects, reconstruction of conditional controls at zero histories, native
full-grade preservation, recursive feedback, or a polynomial substitution
compiler. Those require additional constructions and compatibility checks.
In a control projection whose sole additional datum is a coupling $q$ on
the prescribed Cartesian raw carrier, with all coefficients and native arrow
actions already fixed, this proves uniqueness of a universal native-natural
section whenever such a section exists. It does not prove existence of its
coefficient arrows. Arbitrary joint systems with extra latent witnesses,
different coefficient objects or additional global grades are not classified.
The scalar product law satisfies the usual permutation, unit and grouping
identities, but these scalar equations do not certify corresponding native
coefficient or global-grade coherence.

\begin{thebibliography}{1}
\bibitem{BF} Victor Bloch and Tobias Fritz.
\emph{Naturality characterizes product measures and relative product measures}.
Version 2, 8 October 2026. Theorems 3.1--3.2.
\url{https://arxiv.org/abs/2609.15909v2}.
Primary theorem statements checked at
\url{https://arxiv.org/html/2609.15909v2}.
\end{thebibliography}
\end{document}
